\section{Emergence of Machine Learning in Materials Science} \label{section:machine_learning}

Over the past decade, machine learning (ML) has emerged as a transformative tool in materials modelling, offering a route to accelerate predictions and expand accessible system sizes beyond those tractable by first-principles methods. Traditional computational techniques, such as Density Functional Theory (DFT), Molecular Dynamics (MD), and Kinetic Monte Carlo (KMC), have established themselves as accurate and rigorous, but they often scale poorly with system size or configurational complexity. This limitation becomes particularly acute in the modelling of interfaces, where subtle geometric, chemical, and electronic variations can dramatically affect material behaviour. ML offers a complementary paradigm, enabling the construction of surrogate models trained on high-fidelity data, which can interpolate or even extrapolate to predict key material properties across vast configuration spaces \cite{Leitherer2023-ew, Anderson1960-zp, Low2025-on, Kao2025-ay, Nykiel2023-jf, Gerber2023-yd, Owen2024-wp, Rosen2022-pw, Di_Liberto2022-gz, Chagarov2015-ol, Chen2025-bo, Pyzer-Knapp2025-uq, Davies2021-zh, Taylor2023-wh, Perdew1996-no, Behler2007-nu, Pitfield2025-ah}.

This section introduces the growing role of ML in materials prediction, highlighting its application in interface modelling and the broader landscape of data-driven approaches. Subsequent subsections will examine the types of models in use, including classical regressors and deep neural architectures, alongside the nature of training data, typical input descriptors, and validation strategies. Special attention will be given to machine-learned interatomic potentials (MLPs), particularly those with demonstrated success in interfacial contexts such as MACE and ChgNet \cite{Batatia2022-xz, Maji2021-kc, Davies2024-ju, Cervantes2025-ze, Anderson1960-zp, Batzner2022-fd, Low2025-on, Kao2025-ay, Uhrin2025-ba, Gerber2023-yd, Owen2024-wp, Rosen2022-pw, Di_Liberto2022-gz, Butler2019-ws, Chen2025-bo, Davies2021-zh, Taylor2023-wh, Perdew1996-no, Behler2007-nu, Turlo2019-pd, Pitfield2025-ah}. Finally, we reflect on the benefits and ongoing limitations of ML methods in materials science, including issues of transferability, uncertainty quantification, and integration into automated modelling pipelines \cite{Leitherer2023-ew, Owen2024-wp, Rosen2022-pw, Di_Liberto2022-gz, Butler2019-ws, Chagarov2015-ol, Pyzer-Knapp2025-uq, Taylor2023-wh, Perdew1996-no}.

\subsection{How has machine learning been introduced into materials prediction?}

Machine learning (ML) has been introduced into materials prediction as a complementary tool to traditional first-principles methods, offering improved scalability, predictive flexibility, and the potential for workflow automation. Its adoption has been driven by limitations in conventional approaches such as Density Functional Theory (DFT), particularly in capturing the structural and chemical complexity of real-world systems such as interfaces \cite{Leitherer2023-ew, Low2025-on, Kao2025-ay, Nykiel2023-jf, Gerber2023-yd, Owen2024-wp, Rosen2022-pw, Di_Liberto2022-gz, Butler2019-ws, Davies2021-zh, Taylor2023-wh, Perdew1996-no, Pitfield2025-ah}.

\subsubsection{Computational Motivation}

First-principles methods scale poorly with system size, typically as \( \mathcal{O}(N^3) \), where \( N \) is the number of electrons or basis functions. This makes simulations of large or disordered systems, especially those involving interfaces with long-range strain relaxation, misfit dislocations, or extended defects, computationally prohibitive. Machine-learned interatomic potentials (MLPs) such as MACE or ChgNet offer near-DFT accuracy at a fraction of the cost, enabling simulations with thousands of atoms \cite{Batatia2022-xz, Maji2021-kc, Davies2024-ju, Cervantes2025-ze, Anderson1960-zp, Batzner2022-fd, Low2025-on, Kao2025-ay, Nykiel2023-jf, Uhrin2025-ba, Owen2024-wp, Rosen2022-pw, Di_Liberto2022-gz, Butler2019-ws, Chen2025-bo, Davies2021-zh, Taylor2023-wh, Perdew1996-no, Turlo2019-pd, Pitfield2025-ah}.

\subsubsection{Exploration of Complex Configuration Spaces}

ML enables efficient exploration of vast configurational landscapes, where brute-force DFT sampling would be infeasible. This is particularly important in interface modelling, where multiple stacking orders, terminations, and slip configurations can exist. ML-driven screening allows researchers to rapidly identify promising candidates for further evaluation \cite{Leitherer2023-ew, Low2025-on, Kao2025-ay, Gerber2023-yd, Owen2024-wp, Rosen2022-pw, Di_Liberto2022-gz, Butler2019-ws, Davies2021-zh, Taylor2023-wh, Perdew1996-no}.

\subsubsection{Initial Integration Pathways}

Early ML applications focused on surrogate models trained to replicate DFT-derived energies, forces, or properties. Descriptor-based regression models, using quantities such as coordination numbers, Bader charges, or bond lengths, were employed to predict interfacial stability or band alignments. These approaches have evolved toward end-to-end learning frameworks using physically-informed descriptors or graph-based representations \cite{Leitherer2023-ew, Low2025-on, Kao2025-ay, Owen2024-wp, Rosen2022-pw, Di_Liberto2022-gz, Butler2019-ws, Davies2021-zh, Taylor2023-wh, Perdew1996-no}.

\subsubsection{Workflow Integration}

ML is now embedded into iterative prediction--generation--evaluation cycles. For example, predicted interface favourability informs structure generation (e.g. stacking vectors or surface terminations), which are then evaluated by DFT or ML surrogates. The results feed back into model training, supporting a self-refining loop. Tools such as ARTEMIS benefit from such integration \cite{Leitherer2023-ew, Low2025-on, Kao2025-ay, Gerber2023-yd, Owen2024-wp, Rosen2022-pw, Di_Liberto2022-gz, Butler2019-ws, Davies2021-zh, Taylor2023-wh, Perdew1996-no}.

\subsubsection{Modern Model Types}

More recent work employs graph neural networks (GNNs), such as MACE or NequIP, which are equivariant to Euclidean transformations and thus well-suited for atomistic systems. These models learn directly from atomic positions and species, avoiding the need for handcrafted features while preserving symmetry constraints essential for interface modelling \cite{Leitherer2023-ew, Batzner2022-fd, Low2025-on, Kao2025-ay, Batatia2022-xz, Uhrin2025-ba, Owen2024-wp, Rosen2022-pw, Di_Liberto2022-gz, Butler2019-ws, Davies2021-zh, Taylor2023-wh, Perdew1996-no}.

\subsubsection{Toward Autonomy}

Ultimately, ML enables a shift from static modelling pipelines to more autonomous workflows. When coupled with structure-generation protocols and physical constraints, ML can be used not only to accelerate predictions but to guide the discovery of novel interfacial phenomena in mixed-dimensional and disordered materials \cite{Leitherer2023-ew, Low2025-on, Kao2025-ay, Owen2024-wp, Rosen2022-pw, Di_Liberto2022-gz, Butler2019-ws, Davies2021-zh, Taylor2023-wh, Perdew1996-no}.

\subsection{What types of models and data are being used?}

Recent advances in machine learning (ML) have introduced a diverse range of models for predicting materials properties, including those relevant to interface stability, energy and structure. These models can be broadly categorised into classical learning algorithms and deep learning approaches. Classical methods include random forests, gradient boosting, support vector machines (SVMs), and Gaussian process regression (GPR), the latter being particularly valuable for small datasets due to its uncertainty estimation capability. Deep learning architectures include feedforward neural networks and, more notably, graph neural networks (GNNs), which are well-suited for materials modelling due to their ability to represent variable-sized atomic graphs with complex bonding topologies \cite{Leitherer2023-ew, Low2025-on, Owen2024-wp, Rosen2022-pw, Di_Liberto2022-gz, Butler2019-ws, Taylor2023-wh, Perdew1996-no, Behler2007-nu}.

The input to these models typically consists of numerical representations of atomic configurations known as descriptors. Geometric descriptors include bond lengths, coordination numbers, angular distributions, $n$-body distribution functions, and radial symmetry functions; many of which feature prominently in Behler-Parrinello-type neural networks. Chemical descriptors such as electronegativity differences and Bader charges capture bonding character and charge transfer. Electronic descriptors like partial density of states (PDOS) and local work functions are relevant for predicting band alignment and interfacial transport properties. Interface-specific features, such as interfacial roughness or vacuum level offsets, are particularly crucial in 2D|3D and mixed-dimensional heterostructures \cite{Leitherer2023-ew, Low2025-on, Owen2024-wp, Rosen2022-pw, Di_Liberto2022-gz, Butler2019-ws, Davies2021-zh, Taylor2023-wh, Perdew1996-no}.

To support learning across such heterogeneous and high-dimensional features, symmetry-invariant and scale-independent representations have become increasingly important. These ensure that models generalise across materials with different atom counts or lattice sizes. Dimensionality reduction methods and learned embeddings (e.g. via autoencoders or message passing) are frequently used to compress structural information into a lower-dimensional latent space optimised for prediction \cite{Leitherer2023-ew, Owen2024-wp, Rosen2022-pw, Di_Liberto2022-gz, Butler2019-ws, Taylor2023-wh, Perdew1996-no}.

The datasets used to train such models vary in size and fidelity. They often derive from high-throughput DFT calculations available in public repositories such as the Materials Project or AFLOW. In cases targeting interfaces, the training data may originate from internal simulations using structure generation pipelines like ARTEMIS or RAFFLE, where candidate structures are pre-screened by DFT or empirical methods for energetic viability \cite{Leitherer2023-ew, Nykiel2023-jf, Gerber2023-yd, Owen2024-wp, Rosen2022-pw, Di_Liberto2022-gz, Butler2019-ws, Davies2021-zh, Taylor2023-wh, Perdew1996-no}.

The combination of curated descriptors, structured data from known interfaces, and model architectures sensitive to local and long-range atomic interactions underpins the growing success of ML in materials interface prediction \cite{Leitherer2023-ew, Low2025-on, Owen2024-wp, Rosen2022-pw, Di_Liberto2022-gz, Butler2019-ws, Taylor2023-wh, Perdew1996-no}.

\subsection{What are the strengths and limitations of ML in this field?}

Machine learning (ML) has emerged as a promising strategy for accelerating materials modelling, particularly in the context of complex interfacial systems. Its strengths lie chiefly in efficiency, scalability, and adaptability. ML models, once trained, can rapidly approximate formation energies, surface reactivities, and stability metrics for thousands of configurations, many of which would be prohibitively expensive to assess using first-principles methods alone. This computational efficiency enables exploration of broader configurational and chemical spaces, including systems with disorder, stacking faults, or metastable structures that pose challenges for traditional simulations \cite{Leitherer2023-ew, Low2025-on, Nykiel2023-jf, Gerber2023-yd, Owen2024-wp, Rosen2022-pw, Di_Liberto2022-gz, Butler2019-ws, Davies2021-zh, Taylor2023-wh}.

Another key advantage is the capacity of ML to capture subtle, non-linear correlations in high-dimensional datasets. When used in conjunction with tools such as ARTEMIS or RAFFLE, ML can assist not only in predicting energetics, but also in prioritising candidate geometries during structure generation. Moreover, deep learning approaches such as graph neural networks (GNNs) and equivariant message-passing architectures have demonstrated good performance on systems with varied local coordination environments, as frequently found at 2D|3D or lateral heterointerfaces \cite{Leitherer2023-ew, Low2025-on, Batatia2022-xz, Uhrin2025-ba, Owen2024-wp, Rosen2022-pw, Di_Liberto2022-gz, Butler2019-ws, Taylor2023-wh}.

However, ML approaches also exhibit critical limitations. A foremost concern is their dependence on the training dataset: models typically interpolate well within the domain of known examples but generalise poorly to unseen chemistries or extreme geometries. This limits their reliability when extrapolating to novel material classes or interface types, unless training sets are curated with care to ensure coverage of relevant chemical and structural motifs \cite{Leitherer2023-ew, Low2025-on, Owen2024-wp, Rosen2022-pw, Di_Liberto2022-gz, Butler2019-ws, Taylor2023-wh}.

Another limitation is interpretability. While classical models (e.g. those based on density functional theory) yield physically grounded outputs, such as charge densities or band structures, many ML models operate as statistical black boxes. Though efforts are underway to integrate physically meaningful descriptors (e.g. n-body distribution functions, local work function, or Bader charges), ensuring transparency and robustness remains a work in progress \cite{Leitherer2023-ew, Low2025-on, Owen2024-wp, Rosen2022-pw, Di_Liberto2022-gz, Butler2019-ws, Taylor2023-wh}.

Finally, ML alone cannot fully replace physics-based methods. For systems exhibiting emergent interfacial phases, electronic reconstructions, or coupled phonon-electron effects, high-fidelity methods like DFT or many-body perturbation theory remain necessary for validation and refinement. As such, a hybrid strategy is often advocated; using ML to pre-screen candidate structures and flag configurations for further high-accuracy assessment \cite{Leitherer2023-ew, Hepplestone2006-lo, Low2025-on, Gerber2023-yd, Owen2024-wp, Rosen2022-pw, Di_Liberto2022-gz, Butler2019-ws, Taylor2023-wh}.

\paragraph{In summary} ML constitutes a valuable addition to the materials modelling toolkit, particularly for screening and guiding interface prediction. Yet its performance remains contingent on data quality and physical context, and it is best deployed in tandem with more rigorous electronic structure methods \cite{Leitherer2023-ew, Low2025-on, Owen2024-wp, Rosen2022-pw, Di_Liberto2022-gz, Butler2019-ws, Taylor2023-wh}.

\subsection{Machine-Learned Interatomic Potentials (MLPs)}

Machine-learned interatomic potentials (MLPs) have emerged as a powerful alternative to classical force fields and first-principles methods for atomistic simulations. By learning from high-fidelity data, typically generated from Density Functional Theory (DFT), these models aim to reproduce the accuracy of ab initio methods while offering orders-of-magnitude reductions in computational cost. This is particularly valuable for interface modelling, where large supercells and long-range relaxation effects often render direct DFT approaches intractable \cite{Leitherer2023-ew, Low2025-on, Nykiel2023-jf, Owen2024-wp, Rosen2022-pw, Di_Liberto2022-gz, Butler2019-ws, Davies2021-zh, Taylor2023-wh}.

\subsubsection{MACE}

The MACE (Many-body Attention and Correlation Equivariant) model represents a recent advancement in equivariant message passing neural networks (MPNNs), tailored for high-accuracy and scalable force field construction. Unlike earlier MPNNs which rely primarily on two-body invariant messages, MACE employs many-body interactions and equivariant tensor contractions to capture angular and chemical complexity with improved expressivity and data efficiency \cite{Batatia2022-xz, Uhrin2025-ba, Owen2024-wp, Rosen2022-pw, Di_Liberto2022-gz, Butler2019-ws, Davies2021-zh}.

One notable strength of MACE is its capacity to reach state-of-the-art accuracy on diverse benchmarks, including rMD17, 3BPA, and AcAc, while requiring only two message passing iterations, due to its use of high-body-order messages. For instance, in extrapolation tasks involving out-of-distribution molecular geometries or high-temperature trajectories, MACE (particularly with $L=2$) has demonstrated superior performance and robustness relative to competing models such as NequIP and BOTNet \cite{Batatia2022-xz, Owen2024-wp, Rosen2022-pw, Di_Liberto2022-gz, Butler2019-ws, Davies2021-zh}.

Critically for materials modelling, MACE supports parallel evaluation across many GPUs by limiting its receptive field through local message construction and efficient tensor operations. This makes it attractive for large-scale interface studies involving thousands of atoms; especially where relaxation and strain distribution across the interface play a pivotal role \cite{Leitherer2023-ew, Anderson1960-zp, Hepplestone2006-lo, Batzner2022-fd, Lorentzon2025-si, Low2025-on, Kao2025-ay, Batatia2022-xz, Gerber2023-yd, Owen2024-wp, Rosen2022-pw, Di_Liberto2022-gz, Butler2019-ws, Davies2021-zh, Ceperley1980-bs, Turlo2019-pd}.

\subsubsection{ChgNet}

ChgNet is a neural network potential that incorporates both atomic positions and charge distributions into its prediction framework, thereby explicitly learning charge transfer phenomena. This is particularly relevant for interfaces where electrostatics, dipole formation, and local charge accumulation significantly affect electronic properties such as band alignment.

While detailed benchmark results for ChgNet are less mature compared to MACE, its underlying design philosophy, embedding local charge environments directly into the model, offers a complementary advantage. For instance, ChgNet could be better suited to capturing interface-induced dipoles, charge spillover, or mixed-valence behaviour, which are challenging to encode via purely structural descriptors \cite{Leitherer2023-ew, Low2025-on, Owen2024-wp, Rosen2022-pw, Di_Liberto2022-gz, Butler2019-ws, Pitfield2025-ah}.

\subsubsection{Use in Interface Modelling}

In the context of this project, MLPs like MACE and ChgNet offer distinct but synergistic advantages. MACE, with its speed, high accuracy, and symmetry-respecting architecture, is well-suited for structure relaxation and force prediction across multi-material junctions. ChgNet, on the other hand, may provide more faithful predictions of electronic interface effects arising from local charge redistribution.

Nonetheless, these models are not without limitations. MLPs depend heavily on the representativeness of their training data and can struggle to generalise to unseen chemistries or dissimilar interface environments. Therefore, hybrid workflows that use MLPs for broad structural exploration, followed by targeted DFT recalculations for ambiguous configurations, are increasingly being adopted \cite{Leitherer2023-ew, Lorentzon2025-si, Low2025-on, Gerber2023-yd, Owen2024-wp, Rosen2022-pw, Di_Liberto2022-gz, Butler2019-ws, Davies2021-zh, Pitfield2025-ah}.

The eventual goal is to integrate MLPs into an active learning pipeline alongside tools like RAFFLE or ARTEMIS, facilitating the generation, screening, and optimisation of candidate interface structures at scale \cite{Leitherer2023-ew, Low2025-on, Gerber2023-yd, Owen2024-wp, Rosen2022-pw, Di_Liberto2022-gz, Butler2019-ws, Davies2021-zh, Pitfield2025-ah}.